% Quarto 會自行產生 \documentclass；此檔不可再宣告文件類別。
% fontspec 與 xeCJK 由 Quarto/XeLaTeX 字型設定載入。

\usepackage{amsmath,amssymb,amsthm}
\usepackage{enumerate}
\usepackage{graphicx}
\usepackage{subfig}
\usepackage{float}
\usepackage{wrapfig}
\usepackage[outercaption]{sidecap}
\usepackage{array}
\usepackage{booktabs}
\usepackage{xcolor}
\usepackage{longtable}
\usepackage{colortbl}
\usepackage{listings}
\usepackage[parfill]{parskip}

% 額外英文字型命令
\newfontfamily{\E}{Arial}
\newfontfamily{\A}{Arial}
\newfontfamily{\C}[Scale=0.9]{Arial}
\newfontfamily{\TT}[Scale=0.8]{Times New Roman}
\newfontfamily{\SA}[Scale=0.9]{Arial}
\newfontfamily{\SCP}{Arial}

% 額外中文字型命令（使用 Windows 字型的內部英文名稱較穩定）
\newCJKfontfamily{\MB}{Microsoft JhengHei}
\newCJKfontfamily{\SM}[Scale=0.8]{PMingLiU}
\newCJKfontfamily{\K}{DFKai-SB}

% 中文自動換行
\XeTeXlinebreaklocale "zh"
\XeTeXlinebreakskip=0pt plus 1pt

% 自訂命令
\newcommand{\cw}{\texttt{cw}\kern-.6pt\TeX}
\newcommand{\imgdir}{../PyImages/}

% 中文圖表名稱
\renewcommand{\tablename}{表}
\renewcommand{\figurename}{圖}

% 定理環境
\theoremstyle{remark}
\newtheorem{example}{\MB 範例}
\newtheorem{proj}{\LARGE\MB 專題}

% 顏色
\definecolor{slight}{gray}{0.9}
\definecolor{airforceblue}{rgb}{0.36,0.54,0.66}
\definecolor{arylideyellow}{rgb}{0.91,0.84,0.42}
\definecolor{babyblue}{rgb}{0.54,0.81,0.94}
\definecolor{cadmiumred}{rgb}{0.89,0.0,0.13}
\definecolor{coolblack}{rgb}{0.0,0.18,0.39}
\definecolor{beaublue}{rgb}{0.74,0.83,0.9}
\definecolor{beige}{rgb}{0.96,0.96,0.86}
\definecolor{bisque}{rgb}{1.0,0.89,0.77}
\definecolor{grayxeleven}{rgb}{0.75,0.75,0.75}
\definecolor{limegreen}{rgb}{0.2,0.8,0.2}
\definecolor{splashedwhite}{rgb}{1.0,0.99,1.0}

% lstlisting 程式碼樣式
\lstset{
  language=Python,
  frame=lines,
  basicstyle=\SCP\normalsize,
  keywordstyle=\color{blue},
  commentstyle=\color{airforceblue}\itshape,
  backgroundcolor=\color{beige},
  breaklines=true,
  showstringspaces=false,
  tabsize=2
}
